\documentclass[11pt]{article}

\usepackage[T1]{fontenc}
\usepackage[utf8]{inputenc}
\usepackage{lmodern}
\usepackage{amsmath,amssymb,amsfonts,amsthm}
\usepackage[margin=1in]{geometry}
\usepackage{microtype}
\usepackage{xcolor}
\usepackage[colorlinks=true,linkcolor=teal,urlcolor=teal]{hyperref}

\newcommand{\R}{\mathbb R}
\newcommand{\dd}{\,\mathrm d}
\newcommand{\B}{\mathrm B}
\newcommand{\Bcl}{B_{\mathrm{cl}}}
\newcommand{\diag}{\operatorname{diag}}
\newcommand{\htile}{\mathsf H}

\newtheorem*{lemmaBone}{Lemma B.1 (Uniform Young local limit, rephrased)}
\newtheorem*{lemmaBtwo}{Lemma B.2 (Continuous limit of randomized dilation,
rephrased)}
\newtheorem*{propB}{Randomized Young-label dilation (rephrased)}

\title{A Short Guide to Appendix B\\
\large Young local limits and randomized rounding}
\author{}
\date{}

\begin{document}
\maketitle

\begin{abstract}
Appendix B proves that the Schur--Weyl Young-label law has a uniform Gaussian
local limit and that randomized dilation of a label produces the correctly
dilated Gaussian limit.  The first lemma compares the Young law with a
multinomial law and then proves a cellwise local central limit theorem.  The
second lemma shows that the randomized rounding kernel becomes ordinary
Gaussian dilation in central-limit coordinates.  A final typical-set
argument proves that the fallback correction is asymptotically irrelevant
and that all typical input--output labels satisfy the required branching
condition.  This note explains the geometry of the lattice cells, both
lemmas, and the final proposition.
\end{abstract}

\section{Purpose of Appendix B}

Let
\[
  \Delta_d^\circ
  =\left\{p\in\R^d:p_1>\cdots>p_d>0,
    \ \sum_i p_i=1\right\},
\]
and set $\rho_p=\diag(p_1,\ldots,p_d)$.  Under the Schur--Weyl decomposition
\[
  (\mathbb C^d)^{\otimes N}
  \cong\bigoplus_{\lambda\in\mathsf Y_N}
  \mathcal H_\lambda\otimes\mathcal S_\lambda,
\]
let $\Pi_\lambda$ denote the projection onto the summand indexed by
$\lambda$.  The Young-label probability is
\[
  P_{N,p}(\lambda)
  :=\Pr_{\rho_p^{\otimes N}}\{\text{label }\lambda\}
  =\operatorname{Tr}(\Pi_\lambda\rho_p^{\otimes N})
  =\dim\mathcal S_\lambda\,s_\lambda(p),
  \qquad \lambda\in\mathsf Y_N.
\]
Here $\mathsf Y_N$ is the set of partitions of $N$ with at most $d$ rows,
$\mathcal H_\lambda$ is the corresponding irreducible $\mathrm U(d)$
module, $\mathcal S_\lambda$ is the Specht module, and
$s_\lambda(p)=\operatorname{Tr}_{\mathcal H_\lambda}
\pi_\lambda(\rho_p)$ is the Schur polynomial evaluated at $p$.
More precisely, if $\lambda\sim P_{N,p}$, then
$(\lambda-Np)/\sqrt N$ converges to the centered Gaussian
$\mathsf N_{0,\Sigma_p}$ on $\mathsf H$, where
$\Sigma_p=(\diag(p)-pp^{\mathsf T})|_{\mathsf H}$; Lemma B.1 strengthens
this to uniform $L^1$ convergence after cell embedding.
With respect to the induced Lebesgue measure on $\mathsf H$,
\[
  \mathsf N_{0,\Sigma_p}(\dd x)
  =\varphi_{0,\Sigma_p}(x)\dd x
  =\frac{\exp\!\left[-\tfrac12
    \langle x,\Sigma_p^{-1}x\rangle\right]}
  {(2\pi)^{(d-1)/2}\sqrt{\det_{\mathsf H}\Sigma_p}}\dd x.
\]

For unknown spectrum, the channel cannot use $p$ to sample a fresh target
label.  Instead, starting from an observed input label $\mu\in\mathsf Y_n$,
it dilates $\mu$ by
\[
  \gamma_n=\frac{m_n}{n}\longrightarrow\gamma>1
\]
and randomly rounds the result to the target lattice.  Appendix B proves two
facts needed for this construction:
\begin{enumerate}
  \item $P_{N,p}$ has the same uniform local Gaussian limit as the
  multinomial law with parameter $p$;
  \item randomized rounding sends this limit to the Gaussian with covariance
  $\gamma\Sigma_p$ and has asymptotic affinity $\Bcl(\gamma,d)$ with the
  target law $P_{m_n,p}$.
\end{enumerate}
The last part of the appendix verifies that, on the typical set, the rounded
label is a valid partition and satisfies the representation-theoretic
branching condition.

\section{The affine lattice and its fundamental cell}

Put
\[
  \htile=\left\{x\in\R^d:\sum_i x_i=0\right\},
  \qquad
  \mathsf L_N=\left\{x\in\mathbb Z^d:\sum_i x_i=N\right\}.
\]
The lattice parallel to $\htile$ is generated by
$b_i=e_i-e_d$, $1\leq i<d$.  Appendix B uses the half-open fundamental cell
\begin{equation}\label{eq:cell}
  \mathsf C
  =\left\{\sum_{i=1}^{d-1}t_i b_i:
     -\frac12\leq t_i<\frac12\right\}\subset\htile.
\end{equation}
Write $\mathsf A_N=\{x\in\R^d:\sum_i x_i=N\}$ and let
\[
  \mathcal B_N
  =\bigcup_{\lambda\in\mathsf L_N}\partial(\lambda+\mathsf C).
\]
This countable union of $(d-2)$-dimensional faces has induced
$(d-1)$-dimensional volume $|\mathcal B_N|_{\mathsf A_N}=0$, and
\[
  \mathsf A_N\setminus\mathcal B_N
  =\bigsqcup_{\lambda\in\mathsf L_N}
   \bigl((\lambda+\mathsf C)\setminus\mathcal B_N\bigr).
\]
Thus every point outside $\mathcal B_N$ lies in exactly one translated cell.

More explicitly, let $T:\R^{d-1}\to\htile$ be
$Tt=\sum_{i<d}t_i b_i$.  By definition, the Gram matrix of these coordinate
vectors is $G=T^{\mathsf T}T$, whose entries are
\[
  G_{ij}
  =\left\langle\frac{\partial T}{\partial t_i},
                   \frac{\partial T}{\partial t_j}\right\rangle
  =\langle b_i,b_j\rangle
  =\langle e_i-e_d,e_j-e_d\rangle
  =\delta_{ij}+1.
\]
Thus $G$ has diagonal entries $2$ and off-diagonal entries $1$.  Equivalently,
\[
  G=I_{d-1}+\mathbf 1\mathbf 1^{\mathsf T},
  \qquad \mathbf 1=(1,\ldots,1)^{\mathsf T}\in\R^{d-1}.
\]
Its eigenvalue along $\mathbf 1$ is $d$, while its eigenvalue on
$\mathbf 1^\perp$ is $1$.  Hence
$\det G=d\cdot1^{d-2}=d$.

To define the volume notation, choose an orthonormal basis
$q_1,\ldots,q_{d-1}$ of $\htile$ and let
$Q:\R^{d-1}\to\htile$ be $Qs=\sum_{k<d}s_kq_k$.  For a measurable
$E\subset\htile$, its induced Euclidean volume is
\begin{equation}\label{eq:induced-volume-definition}
  |E|_{\htile}
  :=|Q^{-1}(E)|_{\R^{d-1}}
  =\int_{\R^{d-1}}\mathbf 1_E(Qs)\dd s_1\cdots\dd s_{d-1}.
\end{equation}
Here $\mathbf 1_E(y)=1$ for $y\in E$ and $\mathbf 1_E(y)=0$ for
$y\notin E$.
This does not depend on the chosen orthonormal basis, since two such bases
are related by an orthogonal matrix, whose absolute determinant is $1$.
Equation \eqref{eq:induced-volume-definition} has no Jacobian factor because
$Q$ is an isometry.  In contrast, the columns $b_i=e_i-e_d$ of $T$ are not
orthonormal.  Since $Q^{-1}=Q^{\mathsf T}$ on $\htile$, for every measurable
$A\subset\R^{d-1}$,
\[
  |T(A)|_{\htile}
  =|Q^{-1}T(A)|_{\R^{d-1}}
  =|\det(Q^{-1}T)|\,|A|_{\R^{d-1}}.
\]
Moreover, $QQ^{\mathsf T}$ is the orthogonal projection onto $\htile$ and
$T(\R^{d-1})=\htile$, so
\[
  (Q^{-1}T)^{\mathsf T}(Q^{-1}T)
  =T^{\mathsf T}QQ^{\mathsf T}T
  =T^{\mathsf T}T=G.
\]
Hence $|\det(Q^{-1}T)|=\sqrt{\det G}$, and therefore
\begin{equation}\label{eq:coordinate-volume}
  |T(A)|_{\htile}
  =\sqrt{\det G}\,|A|_{\R^{d-1}}
  =\sqrt d\int_A\dd t_1\cdots\dd t_{d-1}.
\end{equation}
Taking $A=[-1/2,1/2)^{d-1}$, whose coordinate volume is $1$, gives
$|\mathsf C|_{\htile}=\sqrt d$; below $|\mathsf C|$ denotes this induced
volume.

\subsection{Randomized rounding as cell membership}

Let $U_i$ be independent and uniform on $[-1/2,1/2)$ and set
\[
  U=\sum_{i=1}^{d-1}U_i(e_i-e_d).
\]
By \eqref{eq:coordinate-volume}, $U$ is uniform with respect to induced
volume in the concrete sense that
\[
  \Pr(U\in A)=\frac{|A|_{\htile}}{|\mathsf C|_{\htile}}
  \qquad\text{for every measurable }A\subset\mathsf C.
\]
Given $\mu\in\mathsf Y_n$, define
\[
  \widetilde\nu_i
  =\left\lfloor\gamma_n(\mu_i+U_i)+\frac12\right\rfloor
  \quad(i<d),
  \qquad
  \widetilde\nu_d=m_n-\sum_{i<d}\widetilde\nu_i.
\]
Outside a null set of cell boundaries, this is equivalent to
\begin{equation}\label{eq:cell-membership}
  \gamma_n(\mu+U)\in\widetilde\nu+\mathsf C.
\end{equation}
Indeed, for $i<d$, membership in the corresponding coordinate interval says
that $\widetilde\nu_i$ is the nearest integer to
$\gamma_n(\mu_i+U_i)$.  The last coordinate follows automatically because
both vectors in \eqref{eq:cell-membership} have coordinate sum $m_n$.

Let $D_n(\widetilde\nu\mid\mu)$ be the distribution of the raw rounded
label.  Since $U$ is uniform on $\mathsf C$,
\begin{align}
  D_n(\widetilde\nu\mid\mu)
  &=\frac1{|\mathsf C|}
    \left|\left\{u\in\mathsf C:
    \gamma_n(\mu+u)\in\widetilde\nu+\mathsf C\right\}\right|.
  \label{eq:probability-before-change}
\end{align}
Under the change of variables $x=\gamma_n(\mu+u)$, volume in the
$(d-1)$-dimensional space $\htile$ is multiplied by
$\gamma_n^{d-1}$.  The image of the set in
\eqref{eq:probability-before-change} is exactly
\[
  \gamma_n(\mu+\mathsf C)\cap(\widetilde\nu+\mathsf C).
\]
Therefore
\begin{equation}\label{eq:cell-kernel-note}
  D_n(\widetilde\nu\mid\mu)
  =\frac{
    |\gamma_n(\mu+\mathsf C)\cap(\widetilde\nu+\mathsf C)|}
    {\gamma_n^{d-1}|\mathsf C|}.
\end{equation}
Thus $D_n(\widetilde\nu\mid\mu)$ is the fraction of the uniformly
dilated input cell lying in the target cell indexed by $\widetilde\nu$.
The denominator is the volume of $\gamma_n(\mu+\mathsf C)$.

\section{Embedding lattice laws as densities}

For $p\in\Delta_d^\circ$ and $\lambda\in\mathsf L_N$, define the rescaled
cell
\[
  \mathsf C_{N,p}(\lambda)
  =\frac{\lambda-Np+\mathsf C}{\sqrt N}\subset\htile.
\]
Translation by $\lambda-Np$ does not change volume, while multiplication by
$N^{-1/2}$ in the $(d-1)$-dimensional space $\htile$ multiplies volume by
$(N^{-1/2})^{d-1}$.  Hence
\[
  |\mathsf C_{N,p}(\lambda)|
  =(N^{-1/2})^{d-1}|\mathsf C|
  =\frac{|\mathsf C|}{N^{(d-1)/2}}.
\]
These cells tile $\htile$ up to null boundaries.  Let
$R:\mathsf L_N\to\R$ satisfy
$\sum_{\lambda\in\mathsf L_N}|R(\lambda)|<\infty$; this is a finite signed
measure on the discrete set $\mathsf L_N$.  If additionally $R(\lambda)\geq0$
and $\sum_\lambda R(\lambda)=1$, then $R$ is a probability law.  Define
\begin{equation}\label{eq:embedding}
  (\mathcal E_{N,p}R)(x)
  =\frac{N^{(d-1)/2}}{|\mathsf C|}R(\lambda),
  \qquad x\in\mathsf C_{N,p}(\lambda).
\end{equation}
The factor in \eqref{eq:embedding} is the reciprocal cell volume.  Hence
\[
  \int_{\mathsf C_{N,p}(\lambda)}
  (\mathcal E_{N,p}R)(x)\dd x=R(\lambda).
\]
It follows cell by cell that
\begin{align}
  \|\mathcal E_{N,p}R-\mathcal E_{N,p}S\|_1
  &=\|R-S\|_{\ell^1},
  \label{eq:l1-isometry}\\
  \int_{\htile}\sqrt{(\mathcal E_{N,p}R)(x)
                    (\mathcal E_{N,p}S)(x)}\dd x
  &=\sum_{\lambda\in\mathsf L_N}\sqrt{R(\lambda)S(\lambda)}.
  \label{eq:affinity-isometry}
\end{align}
Thus the embedding preserves both total variation in the $\ell^1$
convention and Hellinger affinity.

The limiting covariance on $\htile$ is
\[
  \Sigma_p=\diag(p)-pp^{\mathsf T}.
\]
It has kernel spanned by $(1,\ldots,1)$ and is positive definite after
restriction to $\htile$.  For $x\in\htile$,
\begin{equation}\label{eq:quadratic-form}
  \langle x,\Sigma_p^{-1}x\rangle=\sum_i\frac{x_i^2}{p_i},
  \qquad
  \det_{\htile}\Sigma_p=d\prod_i p_i.
\end{equation}
Indeed, for $x\in\htile$, the vector
\[
  z=\diag(p)^{-1}x
    -\frac1d\left(\sum_i\frac{x_i}{p_i}\right)\mathbf 1
\]
satisfies $\sum_i z_i=\sum_i x_i/p_i-\sum_i x_i/p_i=0$, so
$z\in\htile$.  Since $\Sigma_p\mathbf 1=0$ and $\sum_i x_i=0$,
\[
  \Sigma_p z=x-p\sum_i x_i=x,
  \qquad
  \langle x,z\rangle
  =\sum_i\frac{x_i^2}{p_i}
   -\frac1d\left(\sum_i\frac{x_i}{p_i}\right)\sum_i x_i
  =\sum_i\frac{x_i^2}{p_i},
\]
which proves the first identity.  For the second, let
$\Sigma_p^{(k)}$ be the matrix obtained by deleting row and column $k$.
The matrix determinant lemma gives
\[
  \det\Sigma_p^{(k)}
  =\left(\prod_{i\ne k}p_i\right)
    \left(1-\sum_{i\ne k}p_i\right)
  =\prod_i p_i.
\]
Because the normalized kernel vector is $d^{-1/2}\mathbf 1$, orthogonal
diagonalization gives
\[
  \operatorname{adj}(\Sigma_p)
  =\det_{\htile}(\Sigma_p)
    \frac{\mathbf 1\mathbf 1^{\mathsf T}}d,
  \qquad
  \det\Sigma_p^{(k)}
  =\frac1d\det_{\htile}\Sigma_p.
\]
Consequently, $\det_{\htile}\Sigma_p=d\prod_i p_i$.

\section{Lemma B.1: the uniform Young local limit}

\begin{lemmaBone}
For every compact $K\Subset\Delta_d^\circ$,
\begin{equation}\label{eq:young-local-limit-note}
  \sup_{p\in K}
  \|\mathcal E_{N,p}P_{N,p}-\varphi_{0,\Sigma_p}\|_1
  \longrightarrow0.
\end{equation}
\end{lemmaBone}

This is stronger than weak convergence.  It says that after spreading each
lattice mass uniformly over its cell, the entire density converges in
$L^1$, uniformly over $p$ in a compact subset of the strictly ordered
simplex.

\subsection*{Step 1: compare the Young law with a multinomial law}

Let $X_1,\ldots,X_N$ be independent random variables with
$\Pr(X_k=i)=p_i$, and let
$\Lambda_i=|\{k:X_k=i\}|$.  Thus $\Lambda_i$ counts outcomes of type $i$ and
$\sum_i\Lambda_i=N$.  The law $M_{N,p}$ of
$\Lambda=(\Lambda_1,\ldots,\Lambda_d)$ is
\[
  M_{N,p}(\lambda)
  =\begin{cases}
    \displaystyle\frac{N!}{\prod_i\lambda_i!}
      \prod_i p_i^{\lambda_i},
      &\lambda\in\mathbb Z_{\geq0}^d,
       \ \sum_i\lambda_i=N,\\[2mm]
    0,&\lambda\in\mathsf L_N\setminus\mathbb Z_{\geq0}^d.
  \end{cases}
\]
Define the central window
\[
  \mathsf W_{N,p}
  =\left\{\lambda\in\mathsf L_N:
    \|\lambda-Np\|_\infty\leq N^{2/3}\right\}.
\]
Compactness of $K$ gives uniform lower bounds on every $p_i$ and every gap
$p_i-p_{i+1}$.  Consequently, every nonnegative composition in this window
is a partition for all large $N$.

For a partition $\lambda$, the hook-length and Weyl character formulas give
\begin{align*}
  \dim\mathcal S_\lambda
  &=N!\frac{\prod_{i<j}(\lambda_i-\lambda_j+j-i)}
    {\prod_i(\lambda_i+d-i)!},\\
  s_\lambda(p)
  &=\frac{\det[p_i^{\lambda_j+d-j}]_{i,j}}
    {\prod_{i<j}(p_i-p_j)}.
\end{align*}
On $\mathsf W_{N,p}$, both the exponents $\lambda_j+d-j$ and the numbers
$\log p_i$ are uniformly strictly decreasing.  The rearrangement inequality
therefore makes the identity-permutation term in the determinant strictly
larger than every other term.  The gaps are of order $N$, so the ratio of
each nonidentity term to the identity term is $O(e^{-cN})$, uniformly for
$p\in K$.  Hence
\[
  s_\lambda(p)
  =\frac{\prod_i p_i^{\lambda_i+d-i}}
    {\prod_{i<j}(p_i-p_j)}
    \bigl(1+O(e^{-cN})\bigr).
\]

Dividing $P_{N,p}(\lambda)$ by $M_{N,p}(\lambda)$ now gives
\begin{align}
  \frac{P_{N,p}(\lambda)}{M_{N,p}(\lambda)}
  ={}&\bigl(1+O(e^{-cN})\bigr)
  \prod_{i<j}
  \frac{\lambda_i-\lambda_j+j-i}{N(p_i-p_j)}\notag\\
  &\times
  \prod_i
  \frac{\lambda_i!(Np_i)^{d-i}}{(\lambda_i+d-i)!}.
  \label{eq:young-mult-ratio-note}
\end{align}
Every factor in \eqref{eq:young-mult-ratio-note} is
$1+O_K(N^{-1/3})$ on the central window.  For example,
\[
  \frac{\lambda_i!(Np_i)^{d-i}}{(\lambda_i+d-i)!}
  =\prod_{a=1}^{d-i}\frac{Np_i}{\lambda_i+a}
  =1+O_K(N^{-1/3}).
\]
Therefore $P_{N,p}/M_{N,p}\to1$ uniformly on $\mathsf W_{N,p}$.

Outside the window, the uniform Young concentration estimate controls
$P_{N,p}$ and Hoeffding's inequality controls $M_{N,p}$.  Combining the
central comparison with both tail estimates yields
\begin{equation}\label{eq:young-mult-l1-note}
  \sup_{p\in K}\|P_{N,p}-M_{N,p}\|_{\ell^1}\longrightarrow0.
\end{equation}

\subsection*{Step 2: prove a cellwise local limit for the multinomial law}

Choose
\[
  \frac12<\beta_0<\frac23,
  \qquad
  V_{N,p}=\{\lambda\in\mathsf L_N:
  \|\lambda-Np\|_\infty\le N^{\beta_0}\},
  \qquad
  x_\lambda=\frac{\lambda-Np}{\sqrt N},
\]
and let
\[
  U_{N,p}=\bigcup_{\lambda\in V_{N,p}}
  \mathsf C_{N,p}(\lambda).
\]
Uniform Stirling expansion and a third-order Taylor expansion give, whenever
$\lambda\in V_{N,p}$,
\begin{align}
  \log M_{N,p}(\lambda)
  ={}&-\frac{d-1}{2}\log(2\pi N)
      -\frac12\sum_i\log p_i
      -\frac12\sum_i\frac{x_{\lambda,i}^2}{p_i}\notag\\
     &+O_K\!\left(N^{3\beta_0-2}
                  +N^{\beta_0-1}+N^{-1}\right).
  \label{eq:stirling-note}
\end{align}
The restriction $\beta_0<2/3$ makes the remainder tend to zero.  Using
\eqref{eq:quadratic-form}, $|\mathsf C|=\sqrt d$, and
$\det_{\htile}\Sigma_p=d\prod_i p_i$, exponentiating
\eqref{eq:stirling-note} gives
\begin{equation}\label{eq:pointwise-local-note}
  \frac{N^{(d-1)/2}}{|\mathsf C|}M_{N,p}(\lambda)
  =\varphi_{0,\Sigma_p}(x_\lambda)(1+o(1))
\end{equation}
uniformly on $V_{N,p}$.

Equation \eqref{eq:pointwise-local-note} compares the embedded constant on a
cell with the Gaussian at the cell center.  We first make both tail bounds
explicit.  If $x\in\mathsf C_{N,p}(\lambda)$, then
$x=x_\lambda+u/\sqrt N$ for some $u\in\mathsf C$, and
$\|u\|_\infty\le(d-1)/2$.  Hence
\[
  \|x_\lambda\|_\infty
  \le\|x\|_\infty+\frac{d-1}{2\sqrt N},
  \qquad
  \left\{x:\|x\|_\infty\le
  N^{\beta_0-1/2}-\frac{d-1}{2\sqrt N}\right\}
  \subset U_{N,p}.
\]
For the inclusion, take any $x$ in the set on the left and let
$\mathsf C_{N,p}(\lambda)$ be its cell.  The first inequality gives
$\|x_\lambda\|_\infty\le N^{\beta_0-1/2}$, equivalently
$\|\lambda-Np\|_\infty\le N^{\beta_0}$.  Thus
$\lambda\in V_{N,p}$ and this cell is one of those forming $U_{N,p}$.

If $X\sim M_{N,p}$, then $X_i\sim\operatorname{Bin}(N,p_i)$ for each
$i$.  Coordinatewise Hoeffding followed by a union bound gives
\[
\begin{aligned}
  M_{N,p}(V_{N,p}^c)
  &=\Pr\!\left\{\max_i|X_i-Np_i|>N^{\beta_0}\right\}\\
  &\le\sum_{i=1}^d
       \Pr\!\left\{|X_i-Np_i|>N^{\beta_0}\right\}
   \le 2d\exp\!\left(-2N^{2\beta_0-1}\right).
\end{aligned}
\]
For $Z_p\sim\mathsf N_{0,\Sigma_p}$, put
$r_N=N^{\beta_0-1/2}-(d-1)/(2\sqrt N)$.  The preceding inclusion gives
$U_{N,p}^c\subset\{x:\|x\|_\infty>r_N\}$.  Moreover,
$Z_{p,i}$ is centered Gaussian with variance
$(\Sigma_p)_{ii}=p_i(1-p_i)\le\sigma_K^2$, where
$\sigma_K^2:=\sup_{p\in K,i}p_i(1-p_i)<\infty$.  Hence
\[
\begin{aligned}
  \int_{U_{N,p}^c}\varphi_{0,\Sigma_p}(x)\dd x
  &=\Pr\{Z_p\notin U_{N,p}\}
   \le\Pr\!\left\{\|Z_p\|_\infty>r_N\right\}\\
  &\le\sum_{i=1}^d\Pr\{|Z_{p,i}|>r_N\}
   \le 2d\exp\!\left(-\frac{r_N^2}{2\sigma_K^2}\right)
   \le C_K e^{-c_KN^{2\beta_0-1}}.
\end{aligned}
\]
The last inequality uses
$r_N\ge\tfrac12N^{\beta_0-1/2}$ for all sufficiently large $N$.
Thus the same exponent controls both tails.  Finally, a rescaled cell has
diameter $O(N^{-1/2})$.  On $V_{N,p}$,
$\|\nabla\log\varphi_{0,\Sigma_p}(x)\|=O_K(1+\|x\|)$, so the relative
oscillation of the Gaussian across a cell is
$\exp(O_K(N^{-1/2}+N^{\beta_0-1}))-1=o(1)$.
Thus
\begin{equation}\label{eq:mult-gauss-note}
  \sup_{p\in K}
  \|\mathcal E_{N,p}M_{N,p}-\varphi_{0,\Sigma_p}\|_1
  \longrightarrow0.
\end{equation}
Finally, the isometry \eqref{eq:l1-isometry}, together with
\eqref{eq:young-mult-l1-note} and \eqref{eq:mult-gauss-note}, proves
\eqref{eq:young-local-limit-note}.

\section{Lemma B.2: randomized rounding becomes Gaussian dilation}

Let $\widetilde Q_{m_n,p}=P_{n,p}D_n$ be the raw output law: first draw
$\mu\sim P_{n,p}$ and then
$\widetilde\nu\sim D_n(\cdot\mid\mu)$, before sending invalid outcomes to
the fallback label.

\begin{lemmaBtwo}
For every compact $K\Subset\Delta_d^\circ$,
\begin{align}
  \sup_{p\in K}
  \|\mathcal E_{m_n,p}\widetilde Q_{m_n,p}
     -\varphi_{0,\gamma\Sigma_p}\|_1
  &\longrightarrow0,
  \label{eq:rounded-limit-note}\\
  \sup_{p\in K}
  \left|\B(\widetilde Q_{m_n,p},P_{m_n,p})
        -\Bcl(\gamma,d)\right|
  &\longrightarrow0.
  \label{eq:rounded-affinity-note}
\end{align}
\end{lemmaBtwo}

\subsection*{Step 1: identify the exact continuous operation}

For a density $f$ on $\htile$ and $a>0$, define
\[
  (\operatorname{dil}_a f)(y)
  =a^{-(d-1)/2}f(y/\sqrt a).
\]
This is the density of $\sqrt a X$ when $X$ has density $f$, and it is an
$L^1$ isometry.  Define cell averaging by
\[
  (\operatorname{Av}_{m_n,p}f)(y)
  =\frac{1}{|\mathsf C_{m_n,p}(\nu)|}
    \int_{\mathsf C_{m_n,p}(\nu)}f(z)\dd z
  =\frac{m_n^{(d-1)/2}}{|\mathsf C|}
    \int_{\mathsf C_{m_n,p}(\nu)}f(z)\dd z,
  \qquad y\in\mathsf C_{m_n,p}(\nu).
\]
Thus $\operatorname{Av}_{m_n,p}f$ is constant on every target cell.  This
operator is an $L^1$ contraction.

The central-limit coordinates of a point $\mu+u$ in an input cell obey
\begin{equation}\label{eq:coordinate-dilation-note}
  \frac{\gamma_n(\mu+u)-m_np}{\sqrt{m_n}}
  =\sqrt{\gamma_n}\,
    \frac{\mu+u-np}{\sqrt n}.
\end{equation}
The factor is $\sqrt{\gamma_n}$ rather than $\gamma_n$ because the input
and target central-limit scalings use $\sqrt n$ and $\sqrt{m_n}$.
Equation \eqref{eq:cell-kernel-note} then says exactly which target cell
contains the dilated point.  Set
$g_n=\operatorname{dil}_{\gamma_n}(\mathcal E_{n,p}P_{n,p})$.  On the
dilated input cell $\sqrt{\gamma_n}\,\mathsf C_{n,p}(\mu)$, with
$\mu\in\mathsf Y_n$,
\[
  g_n(y)
  =\gamma_n^{-(d-1)/2}(\mathcal E_{n,p}P_{n,p})
    (y/\sqrt{\gamma_n})
  =\frac{n^{(d-1)/2}}{\gamma_n^{(d-1)/2}|\mathsf C|}P_{n,p}(\mu).
\]
Equation \eqref{eq:coordinate-dilation-note} also gives
\[
  \sqrt{\gamma_n}\,\mathsf C_{n,p}(\mu)
  =\frac{\gamma_n(\mu+\mathsf C)-m_np}{\sqrt{m_n}},
\]
and therefore
\[
  \left|\mathsf C_{m_n,p}(\nu)
  \cap\sqrt{\gamma_n}\,\mathsf C_{n,p}(\mu)\right|
  =\frac{|(\nu+\mathsf C)\cap\gamma_n(\mu+\mathsf C)|}
         {m_n^{(d-1)/2}}.
\]
Multiplying the preceding constant density by this intersection volume and
using $m_n=\gamma_n n$ and \eqref{eq:cell-kernel-note} gives
\begin{align*}
  &\int_{\mathsf C_{m_n,p}(\nu)
     \cap\sqrt{\gamma_n}\,\mathsf C_{n,p}(\mu)}g_n(y)\dd y\\
  &\qquad=P_{n,p}(\mu)
    \frac{|(\nu+\mathsf C)\cap\gamma_n(\mu+\mathsf C)|}
         {\gamma_n^{d-1}|\mathsf C|}
  =P_{n,p}(\mu)D_n(\nu\mid\mu).
\end{align*}
Summing over the input cells, for every $\nu\in\mathsf L_{m_n}$,
\[
  \int_{\mathsf C_{m_n,p}(\nu)}g_n(y)\dd y
  =\sum_{\mu\in\mathsf Y_n}P_{n,p}(\mu)D_n(\nu\mid\mu)
  =\widetilde Q_{m_n,p}(\nu).
\]
Hence, for $y\in\mathsf C_{m_n,p}(\nu)$,
\[
  (\operatorname{Av}_{m_n,p}g_n)(y)
  =\frac{\widetilde Q_{m_n,p}(\nu)}
         {|\mathsf C_{m_n,p}(\nu)|}
  =\frac{m_n^{(d-1)/2}}{|\mathsf C|}
    \widetilde Q_{m_n,p}(\nu)
  =(\mathcal E_{m_n,p}\widetilde Q_{m_n,p})(y).
\]
Thus
\begin{equation}\label{eq:intertwining-note}
  \mathcal E_{m_n,p}\widetilde Q_{m_n,p}
  =\operatorname{Av}_{m_n,p}
   \operatorname{dil}_{\gamma_n}
   (\mathcal E_{n,p}P_{n,p}).
\end{equation}
This is an identity, not merely an asymptotic approximation.  Randomizing
within the input cell is what makes the discrete kernel intertwine exactly
with continuous dilation followed by target-cell averaging.

\subsection*{Step 2: pass to the Gaussian limit}

Lemma B.1 and the $L^1$ isometry of dilation give
\[
  \operatorname{dil}_{\gamma_n}
  (\mathcal E_{n,p}P_{n,p})
  -\operatorname{dil}_{\gamma_n}\varphi_{0,\Sigma_p}
  \longrightarrow0
\]
uniformly in $L^1$.  Moreover,
\[
  \operatorname{dil}_{\gamma_n}\varphi_{0,\Sigma_p}
  =\varphi_{0,\gamma_n\Sigma_p}
  \longrightarrow\varphi_{0,\gamma\Sigma_p}
\]
uniformly for $p\in K$.  Finally, the target-cell diameter is
$O(m_n^{-1/2})$.  The compact family of Gaussians
$\varphi_{0,a\Sigma_p}$, with $p\in K$ and $a$ near $\gamma$, is uniformly
translation-continuous in $L^1$, so averaging it over the vanishing cells
changes it by
\[
  \|\operatorname{Av}_{m_n,p}\varphi_{0,\gamma_n\Sigma_p}
    -\varphi_{0,\gamma_n\Sigma_p}\|_1=o(1).
\]
Now \eqref{eq:intertwining-note}, the triangle inequality, and the
contraction of $\operatorname{Av}_{m_n,p}$ give
\[
\begin{aligned}
  &\|\mathcal E_{m_n,p}\widetilde Q_{m_n,p}
       -\varphi_{0,\gamma\Sigma_p}\|_1\\
  &\quad=\|\operatorname{Av}_{m_n,p}\operatorname{dil}_{\gamma_n}
       (\mathcal E_{n,p}P_{n,p})
       -\varphi_{0,\gamma\Sigma_p}\|_1\\
  &\quad\le
       \|\operatorname{Av}_{m_n,p}
       [\operatorname{dil}_{\gamma_n}(\mathcal E_{n,p}P_{n,p})
       -\varphi_{0,\gamma_n\Sigma_p}]\|_1\\
  &\qquad\quad
       +\|\operatorname{Av}_{m_n,p}\varphi_{0,\gamma_n\Sigma_p}
       -\varphi_{0,\gamma_n\Sigma_p}\|_1
       +\|\varphi_{0,\gamma_n\Sigma_p}
       -\varphi_{0,\gamma\Sigma_p}\|_1\\
  &\quad\le
       \|\operatorname{dil}_{\gamma_n}(\mathcal E_{n,p}P_{n,p})
       -\varphi_{0,\gamma_n\Sigma_p}\|_1\\
  &\qquad\quad
       +\|\operatorname{Av}_{m_n,p}\varphi_{0,\gamma_n\Sigma_p}
       -\varphi_{0,\gamma_n\Sigma_p}\|_1
       +\|\varphi_{0,\gamma_n\Sigma_p}
       -\varphi_{0,\gamma\Sigma_p}\|_1=o(1),
\end{aligned}
\]
which proves \eqref{eq:rounded-limit-note}.

\subsection*{Step 3: take the affinity limit}

Lemma B.1 at sample size $m_n$ also gives
\[
  \mathcal E_{m_n,p}P_{m_n,p}
  \longrightarrow\varphi_{0,\Sigma_p}
\]
uniformly in $L^1$.  For probability densities,
\[
  \|\sqrt f-\sqrt g\|_2^2\leq\|f-g\|_1,
\]
so Hellinger affinity is continuous under $L^1$ convergence.  The affinity
is preserved exactly by the cell embedding in
\eqref{eq:affinity-isometry}.  Therefore
\begin{align*}
  \B(\widetilde Q_{m_n,p},P_{m_n,p})
  &\longrightarrow
  \int_{\htile}
  \sqrt{\varphi_{0,\gamma\Sigma_p}(x)
        \varphi_{0,\Sigma_p}(x)}\dd x\\
  &=\left(\frac{2\sqrt\gamma}{1+\gamma}\right)^{(d-1)/2}
   =\Bcl(\gamma,d).
\end{align*}
The covariance $\Sigma_p$ cancels in the centered Gaussian affinity, which
explains why the classical factor depends only on $\gamma$ and $d$.

\section{From the raw kernel to the valid Young-label kernel}

The raw output $\widetilde\nu$ always lies in $\mathsf L_{m_n}$, but it need
not be a partition, and $\widetilde\nu-\mu$ need not be a partition of
$m_n-n$.  The actual kernel $K_n^{\mathrm{rnd}}$ keeps a raw outcome only
when both conditions hold and the outcome is not the fallback label
\[
  \nu_n^{\mathrm{fb}}=(m_n,0,\ldots,0).
\]
All rejected mass is moved to this fallback label.

\begin{propB}
There is $c_{\mathrm{rnd}}<\infty$ such that, uniformly for
$p$ in a compact $K\Subset\Delta_d^\circ$, every sufficiently large $n$,
$\mu$ in the typical set
\[
  \mathsf T_n(p)
  =\left\{\mu\in\mathsf Y_n:
    \|\mu/n-p\|_\infty\leq n^{-1/3}\right\},
\]
and every $\nu$ with $K_n^{\mathrm{rnd}}(\nu\mid\mu)>0$ satisfy
\[
  \nu-\mu\vdash m_n-n,
  \qquad
  \|\nu-\gamma_n\mu\|_\infty\leq c_{\mathrm{rnd}}.
\]
If
\[
  Q_{m_n,p}(\nu)
  =\sum_\mu P_{n,p}(\mu)K_n^{\mathrm{rnd}}(\nu\mid\mu),
\]
then
\[
  \sup_{p\in K}
  |\B(Q_{m_n,p},P_{m_n,p})-\Bcl(\gamma,d)|
  \longrightarrow0.
\]
\end{propB}

\subsection*{Step 1: the rounding error is uniformly bounded}

For $i<d$, nearest-integer rounding gives
\[
  |\widetilde\nu_i-\gamma_n\mu_i|
  \leq\frac{\gamma_n+1}{2}.
\]
The last-coordinate errors sum to the negative of the first $d-1$ errors,
so
\[
  |\widetilde\nu_d-\gamma_n\mu_d|
  \leq\frac{(d-1)(\gamma_n+1)}2.
\]
Since $\gamma_n\to\gamma$, one may take
\begin{equation}\label{eq:crnd-note}
  c_{\mathrm{rnd}}
  =\frac{d-1}{2}\left(1+\sup_n\gamma_n\right)<\infty.
\end{equation}

\subsection*{Step 2: typical labels need no fallback correction}

Compactness of $K$ inside the strictly ordered simplex gives $c_0>0$ such
that, for all large $n$ and every $\mu\in\mathsf T_n(p)$,
\[
  \mu_d\geq c_0n,
  \qquad
  \mu_i-\mu_{i+1}\geq c_0n\quad(i<d).
\]
Set $\delta=\widetilde\nu-\mu$.  It is an integer vector with coordinate
sum $m_n-n$.  The rounding bound gives
\begin{align*}
  \delta_d
  &\geq(\gamma_n-1)\mu_d-c_{\mathrm{rnd}},\\
  \delta_i-\delta_{i+1}
  &\geq(\gamma_n-1)(\mu_i-\mu_{i+1})-2c_{\mathrm{rnd}}.
\end{align*}
Because $\gamma_n-1$ is bounded away from zero for large $n$, both
right-hand sides are positive.  Hence
\[
  \delta_1\geq\cdots\geq\delta_d>0,
\]
so $\delta\vdash m_n-n$.  Since $\mu$ and $\delta$ are partitions,
$\widetilde\nu=\mu+\delta$ is also a partition.  Moreover,
$\widetilde\nu_d>\mu_d>0$, whereas the fallback label has last coordinate
zero.  Thus $D_n(\cdot\mid\mu)$ is supported entirely on accepted labels
and assigns no mass to the fallback label.  For such typical $\mu$,
\[
  K_n^{\mathrm{rnd}}(\nu\mid\mu)=D_n(\nu\mid\mu).
\]

\subsection*{Step 3: replacement changes only negligible mass}

Let $\varepsilon_{n,p}$ be the probability that the raw outcome is replaced
when $\mu\sim P_{n,p}$.  Step 2 shows that replacement can occur only when
$\mu\notin\mathsf T_n(p)$, so
\[
  \varepsilon_{n,p}
  \leq P_{n,p}(\mathsf T_n(p)^c)\longrightarrow0
\]
uniformly on $K$ by uniform Young concentration.

Let $\widetilde Q_{m_n,p}$ be the raw output law and $Q_{m_n,p}$ the law
after replacement.  Moving mass $\varepsilon_{n,p}$ from rejected labels to
one fallback label changes the $\ell^1$ distance by at most twice that mass:
\[
  \|Q_{m_n,p}-\widetilde Q_{m_n,p}\|_{\ell^1}
  \leq2\varepsilon_{n,p}=o(1).
\]
The continuity of Hellinger affinity and Lemma B.2 now give
\[
  \B(Q_{m_n,p},P_{m_n,p})
  =\B(\widetilde Q_{m_n,p},P_{m_n,p})+o(1)
  =\Bcl(\gamma,d)+o(1),
\]
uniformly for $p\in K$.  This proves the proposition.

\section{How the pieces fit together}

The logical flow of Appendix B is
\[
\begin{array}{c}
P_{N,p}\text{ is uniformly close to the multinomial law}\\[1mm]
\Downarrow\\[-1mm]
\mathcal E_{N,p}P_{N,p}\to\varphi_{0,\Sigma_p}
\quad\text{in }L^1
\qquad\text{Lemma B.1}\\[1mm]
\Downarrow\\[-1mm]
\text{randomized rounding}
=\text{ dilation by }\sqrt{\gamma_n}
+\text{ cell averaging}\\[1mm]
\Downarrow\\[-1mm]
\mathcal E_{m_n,p}\widetilde Q_{m_n,p}
\to\varphi_{0,\gamma\Sigma_p}
\qquad\text{Lemma B.2}\\[1mm]
\Downarrow\\[-1mm]
\B(\widetilde Q_{m_n,p},P_{m_n,p})
\to\Bcl(\gamma,d)\\[1mm]
\Downarrow\\[-1mm]
\text{the fallback changes }o(1)\text{ mass}\\[1mm]
\Downarrow\\[-1mm]
\B(Q_{m_n,p},P_{m_n,p})\to\Bcl(\gamma,d).
\end{array}
\]
The randomized cell variable is essential: it turns discrete rounding into
the exact intertwining identity \eqref{eq:intertwining-note}.  The typical-set
argument is logically separate.  It does not establish the Gaussian limit;
it shows that the raw limiting construction is compatible with the Young
branching rule except on an event of vanishing probability.

\section{Points worth keeping straight}

\begin{enumerate}
  \item The raw kernel $D_n$ is defined on the full affine lattice.  The
  corrected kernel $K_n^{\mathrm{rnd}}$ is a Markov kernel on Young labels.
  \item The power $\gamma_n^{d-1}$ in
  \eqref{eq:cell-kernel-note} is a Jacobian in the
  $(d-1)$-dimensional hyperplane $\htile$.
  \item Dilation by $\gamma_n$ in the original labels becomes dilation by
  $\sqrt{\gamma_n}$ after central-limit normalization.
  \item Lemma B.1 is an $L^1$ local limit, not merely convergence in
  distribution.  This strength is what permits passage to Hellinger
  affinity.
  \item Compactness of $K\Subset\Delta_d^\circ$ supplies uniform lower
  bounds for the coordinates $p_i$, the spectral gaps $p_i-p_{i+1}$, and
  the eigenvalues of $\Sigma_p$ on $\htile$.
  \item The fallback label is needed to define a valid kernel for every
  input label, but it is asymptotically invisible because typical labels
  never use it.
\end{enumerate}

\end{document}
